\documentclass[aps,pre,onecolumn,superscriptaddress,floatfix,nofootinbib]{revtex4-2}

\usepackage{amsmath,amssymb,graphicx,hyperref}
\usepackage{longtable}
\usepackage{array}
\newcolumntype{L}[1]{>{\raggedright\arraybackslash}p{#1}}

\newcommand{\ave}[1]{\langle #1 \rangle}
\newcommand{\code}[1]{\texttt{#1}}
\newcommand{\kbar}{\bar\kappa}
\newcommand{\sbar}{\bar s}
\renewcommand{\thesection}{S\arabic{section}}
\renewcommand{\theequation}{S\arabic{equation}}
\renewcommand{\thetable}{S\arabic{table}}

\begin{document}

\title{Supplemental Material:\\ Statistical mechanics on chygraphs}

\author{Alexei V\'azquez}
\email{alexei@nodeslinks.com}
\affiliation{Nodes \& Links Ltd, Salisbury House, Station Road,
Cambridge, CB1 2LA, UK}

\date{\today}

\maketitle

This document carries three things the main text refers to but has no room for.
Section~\ref{sec:map} is the map from every numbered equation of the manuscript
to the routine that evaluates it and the test that checks it.
Sections~\ref{sec:branching}--\ref{sec:core} write out the derivations that the
main text compresses to a sentence, each one closed by hand and each one pinned
by a test rather than by a quoted number --- the manuscript's own record is that
every claim which had to be corrected during review was a number reported
without a derivation behind it. Section~\ref{sec:gbp} reports the
generalised-belief-propagation calculation that the conclusions name as the
standing item, on the manuscript's own example.
Section~\ref{sec:repro} says which script produces each table and figure.

Equation numbers without an ``S'' refer to the main text. All code is at
\url{https://github.com/av2atgh/chygraph_statmech}; a name written
\code{module.thing} is \code{src/chygraph\_statmech/module.py}, and
\code{Chygraph.thing} is a method of the single class in \code{api.py} that
holds the structure and delegates to those modules.

\section{Equation to method}
\label{sec:map}

{\footnotesize
\begin{longtable}{@{}r L{0.29\textwidth} L{0.37\textwidth} L{0.17\textwidth}@{}}
\caption{\label{tab:map}Every numbered equation of the manuscript, the routine
that evaluates it, and the test file that checks it.}\\
\hline\hline
Eq. & what it is & evaluated by & checked in \\
\hline
\endfirsthead
\hline\hline
Eq. & what it is & evaluated by & checked in \\
\hline
\endhead
\hline\hline
\endfoot
(1)  & up step, $h_{i\to a}$ & \code{population.CavityPopulation.sweep} & \code{test\_population} \\
(2)  & down step, $u_{a\to i}$ & \code{cavity.emitted\_field} & \code{test\_ising} \\
(3)  & generating function acting on a measure & \code{population.CavityPopulation.}\newline\code{\_sum\_over\_complexes} & \code{test\_population} \\
(4)  & ensemble map for $P^{ml}$ & \code{population.CavityPopulation.sweep} & \code{test\_population} \\
(5)  & field emitted by a complex & \code{cavity.emitted\_field}, \code{Chygraph.emitted\_field} & \code{test\_ising} \\
(6)  & linearisation, $u'=\partial u/\partial h$ & \code{cavity.cavity\_derivative} & \code{test\_ising} \\
(7)  & reweighting $\ave{\kappa}\to\ave{\kappa}\ave{u'}$ & \code{stability.reweight}, \code{Chygraph.stability\_tensor} & \code{test\_stability} \\
(8)  & branching matrix $B_{lm}$ & \code{Chygraph.branching\_matrix} & \code{test\_api}, \code{test\_derivations} \\
(9)  & size-biased $\ave{\sbar\,u'}_m$ & \code{Chygraph.transmission}, \code{Chygraph.size\_biased} & \code{test\_api}, \code{test\_derivations} \\
(10) & $\det(I-B)=0$ & \code{Chygraph.critical\_coupling} & \code{test\_api}, \code{test\_ising} \\
(11) & single-cardinality collapse of (8) & \code{ising.branching\_matrix} & \code{test\_ising} \\
(12) & $\rho(J(Q^\ast))=1-\Lambda+O(\Lambda^2)$ & \code{fixedpoint.FixedPointStability.}\newline\code{spectral\_radius} & \code{test\_fixedpoint} \\
(13) & Bethe free energy & \code{freeenergy.BetheFreeEnergy.minus\_beta\_f} & \code{test\_freeenergy} \\
(14) & paramagnetic closed form & \code{freeenergy.paramagnetic} & \code{test\_freeenergy}, \code{test\_derivations} \\
(15) & links and triangles, $L=2$ & \code{Chygraph.critical\_coupling} & \code{test\_derivations} \\
(16) & Ising Hamiltonian & \code{cavity.ising\_clique} & \code{test\_ising} \\
(17) & $u'$ for $c=2,3,4$ & \code{ising.clique\_derivative} & \code{test\_ising}, \code{test\_derivations} \\
(18) & Monte Carlo $T_c$ & \code{probe/ising\_mc.py} & --- (simulation) \\
(19) & simplicial Hamiltonian & \code{simplicial.emitted} & \code{test\_simplicial} \\
(20) & $Z(\sigma_0)$ for the unanimity rule & \code{simplicial.emitted} & \code{test\_simplicial}, \code{test\_derivations} \\
(21) & $u'$ for the unanimity rule & \code{simplicial.uprime} & \code{test\_simplicial}, \code{test\_derivations} \\
(22) & Bragg--Williams limit $T^\ast$ & \code{simplicial.SimplicialChygraph.spinodal} & \code{test\_simplicial}, \code{test\_derivations} \\
(23) & first-order $T_c$ and $\Delta m$ at $q=16$ & \code{simplicial.SimplicialChygraph.transition} & \code{test\_simplicial} \\
(24) & hitting-set Hamiltonian and constraint & \code{hittingset.HittingSet} & \code{test\_hittingset} \\
(25) & hard-field warning propagation & \code{hittingset.HittingSet.F} & \code{test\_hittingset} \\
(26) & hard-field ensemble map & \code{hittingset.HittingSet.solve} & \code{test\_hittingset} \\
(27) & $\ave{k}(c-1)=e$ & \code{hittingset.rsb\_point} & \code{test\_hittingset}, \code{test\_derivations} \\
(28) & vertex-cover constraint & \code{cover.CliqueCover} & \code{test\_cover} \\
(29) & the $\max$ rule inside a complex & \code{cover.CliqueCover.tau} & \code{test\_cover} \\
(30) & vertex-cover ensemble map & \code{cover.CliqueCover.solve} & \code{test\_cover} \\
(31) & soft-field belief propagation & \code{softfield.HittingSetBP.sweep} & \code{test\_softfield} \\
(32) & $h_{\mathrm{RS}}$ and $\rho=1/K$ & \code{softfield.regular\_field}, \code{softfield.regular\_density} & \code{test\_softfield}, \code{test\_derivations} \\
(33) & entropy $s$ for an arbitrary chygraph & \code{softfield.HittingSetBP.entropy} & \code{test\_softfield} \\
(34) & regular-hypergraph entropy & \code{softfield.regular\_entropy} & \code{test\_softfield}, \code{test\_derivations} \\
(35) & three-state leaf-removal messages & \code{core.CorePercolation.solve} & \code{test\_core} \\
(36) & $\zeta$ and $\kappa$ inside a complex & \code{core.CorePercolation.zeta}, \code{.kappa} & \code{test\_core}, \code{test\_derivations} \\
(37) & M\"obius counting numbers & \code{region.RegionGraph} & \code{test\_region} \\
\end{longtable}
}

Two entries need a word. Equation~(2) is a placeholder for whatever the
interior Hamiltonian is: the routine takes $-\beta H$ as a callable, and the
four models of the manuscript differ only in what is passed to it.
Equation~(18) is a direct simulation and is deliberately not covered by the test
suite, because a test that reproduced it from the cavity equations would defeat
its purpose; it is the one number in the paper obtained without them.

The remaining routines are solvers shared across sections rather than
equations: \code{antimonotone.fixed\_point} brackets the order-reversing maps of
Eqs.~(26) and~(30) by iterating $F\circ F$ from $0$ and from $1$,
\code{region.overlap\_profile} measures how far an explicit complex family is
from treelike, and \code{gbp.GBP} is Sec.~\ref{sec:gbp} below.

\section{The branching matrix}
\label{sec:branching}

\subsection{The size-biased average, Eq.~(9)}

Arriving at a node through one of its complexes and asking what else branches
means sampling a complex \emph{by following one of its inclusions}. A complex of
cardinality $c$ has $c$ inclusions and so is reached $c$ times as often as one
of cardinality $1$; the ensemble seen on arrival is therefore the size-biased
one, $c\,p_c/\ave{c}$, and every per-complex quantity in Eq.~(8) is averaged
that way. With $f(c)=(c-1)u'(c)$ this is Eq.~(9). Taking instead $f(c)=c-1$,
\begin{equation}
  \ave{\sbar}=\sum_c\frac{c\,p_c}{\ave{c}}(c-1)
  =\frac{\ave{c^{2}}-\ave{c}}{\ave{c}}
  =\frac{\ave{c^{2}}}{\ave{c}}-1 ,
  \label{eq:sbar}
\end{equation}
which is not $\ave{c}-1$: for a layer mixing $c=3$ and $c=5$ equally the two are
$3.25$ and $3$. And $\ave{\sbar\,u'}\neq\ave{\sbar}\ave{u'}$ unless $u'$ is
common to the layer, which is why Eq.~(9) averages the product. Splitting a
mixed layer into one layer per cardinality with chy-degrees in the ratio
$c\,p_c$ reproduces the same transition, which is the check
\code{test\_api} runs.

\subsection{Links and triangles, Eq.~(15)}

With $L$ the link layer ($c=2$) and $T$ the triangle layer ($c=3$), Eq.~(8) reads
\begin{equation}
  B=\begin{pmatrix}
      \ave{\kbar}_L u'_L & 2\ave{\kappa}_T u'_T\\[2pt]
      \ave{\kappa}_L u'_L & 2\ave{\kbar}_T u'_T
    \end{pmatrix},
\end{equation}
the factors $c_m-1$ being $1$ and $2$. Then
\begin{equation}
  \begin{split}
  \det(I-B)&=\big(1-\ave{\kbar}_Lu'_L\big)\big(1-2\ave{\kbar}_Tu'_T\big)
             -2\ave{\kappa}_L\ave{\kappa}_Tu'_Lu'_T\\
  &=1-\ave{\kbar}_Lu'_L-2\ave{\kbar}_Tu'_T
     +2u'_Lu'_T\big(\ave{\kbar}_L\ave{\kbar}_T-\ave{\kappa}_L\ave{\kappa}_T\big),
  \end{split}
\end{equation}
and setting it to zero is Eq.~(15). The cross term carries
$\ave{\kappa}_L\ave{\kappa}_T-\ave{\kbar}_L\ave{\kbar}_T$ and so vanishes
whenever both layers are Poisson, each being its own excess. It does not vanish
for regular layers, where $\ave{\kbar}=\ave{\kappa}-1$, and that is the whole of
the difference between the two null models of Table~I.

\section{The Ising model}
\label{sec:ising}

\subsection{What a clique transmits, Eq.~(17)}

Differentiating Eq.~(5) with respect to one incoming field and evaluating at
$h=0$,
\begin{equation}
  u'=\frac{\partial}{\partial h_1}\,\tfrac12\big[\ln Z(+)-\ln Z(-)\big]\bigg|_{h=0}
   =\tfrac12\Big[\ave{\sigma_1}_+-\ave{\sigma_1}_-\Big],
  \label{eq:updef}
\end{equation}
the two averages being taken inside the complex at zero cavity field with
$\sigma_0$ held at $+1$ and at $-1$. For $c=3$ with $a=\beta J$, the interior
sum $\sigma_0\sigma_1+\sigma_0\sigma_2+\sigma_1\sigma_2$ takes the value $3$ on
the one configuration where the other two agree with $\sigma_0$ and $-1$ on the
other three, so
\begin{equation}
  \ave{\sigma_1}_+=\frac{e^{3a}-e^{-a}}{e^{3a}+3e^{-a}}
                  =\frac{e^{4a}-1}{e^{4a}+3},
  \qquad \ave{\sigma_1}_-=-\ave{\sigma_1}_+
\end{equation}
by spin-flip symmetry. Hence $u'=(e^{4a}-1)/(e^{4a}+3)$, and substituting
$e^{2a}=(1+t)/(1-t)$ with $t=\tanh a$,
\begin{equation}
  u'=\frac{(1+t)^2-(1-t)^2}{(1+t)^2+3(1-t)^2}
    =\frac{4t}{4-4t+4t^{2}}=\frac{t}{1-t+t^{2}} ,
\end{equation}
which is Eq.~(17). The same enumeration at $c=4$ gives
$(t+t^{3})/(1-2t+3t^{2})$, and beyond that no useful closed form, though the sum
remains exact at any cardinality small enough to enumerate. Since
$t/(1-t+t^2)>t$ on $0<t<1$, a triangle transmits more per traversal than an
edge; Sec.~V~B of the main text is the argument that it nevertheless lowers
$T_c$, because Eq.~(8) also loses a branch.

\subsection{Why the overlap term collapses, Eq.~(13)}

The Bethe free energy of a factor model is a sum of one term per factor and one
per variable, minus one per factor--variable pair,
\begin{equation}
  -\beta F=\sum_a\ln Z_a+\sum_i\ln Z_i-\sum_{(a,i)}\ln Z_{ai},
\end{equation}
with $Z_{ai}=\sum_\sigma\nu_{i\to a}(\sigma)\hat\nu_{a\to i}(\sigma)$ the
overlap term. Parameterising $\nu_{i\to a}(\sigma)\propto e^{h_{i\to a}\sigma}$
and $\hat\nu_{a\to i}(\sigma)\propto e^{u_{a\to i}\sigma}$ gives
$Z_{ai}=2\cosh(h_{i\to a}+u_{a\to i})$, and by Eq.~(1)
\begin{equation}
  h_{i\to a}+u_{a\to i}
  =\mu+\sum_{b\ni i,\,b\neq a}u_{b\to i}+u_{a\to i}
  =\mu+\sum_{b\ni i}u_{b\to i}=h_i ,
\end{equation}
the \emph{full} field at $i$, which does not depend on $a$. So every one of the
$k_i$ overlap terms at node $i$ equals $Z_i=2\cosh h_i$, and
\begin{equation}
  \sum_i\ln Z_i-\sum_{(a,i)}\ln Z_{ai}=\sum_i(1-k_i)\ln Z_i .
\end{equation}
Per node, with $n_l=\ave{k_l}/c_l$ layer-$l$ complexes --- each holds $c_l$
members and each node sits in $\ave{k_l}$ of them --- this is Eq.~(13). It is
the collapse that makes the free energy cost no more than the messages already
computed, and it is special to the exponential parameterisation: for a general
message shape the overlap term does not reduce to the node term.

At zero cavity field $Z_i=2$ and $\ln Z_a$ becomes the isolated-complex
$\ln Z_c$, so
\begin{equation}
  -\beta f_{\mathrm{para}}
   =\sum_l n_l\ln Z^{(l)}_c+\Big(1-\sum_l\ave{k_l}\Big)\ln 2
   =\ln 2+\sum_l\frac{\ave{k_l}}{c_l}\ln\!\big(Z^{(l)}_c/2^{c_l}\big),
\end{equation}
which is Eq.~(14). For a graph it is the textbook
$\ln 2+(\ave{k}/2)\ln\cosh\beta J$.

\section{The unanimity interaction}
\label{sec:simplicial}

\subsection{The interior sum, Eqs.~(20) and~(21)}

The unanimity rule contributes $e^{a}$, $a=\beta J_q$, on the two configurations
in which every member of the hyperedge agrees, and $1$ on the other $2^q-2$.
Writing it as $1+(e^{a}-1)\delta$ splits the interior sum into a free part and
one correction term. With independent cavity fields $h_1,\dots,h_{q-1}$ on the
other members,
\begin{equation}
  Z(\sigma_0)=\prod_{j=1}^{q-1}\big(2\cosh h_j\big)
              +(e^{a}-1)\,e^{\sigma_0\sum_j h_j},
  \label{eq:simZgen}
\end{equation}
the correction surviving only where every $\sigma_j=\sigma_0$. Putting every
$h_j=h$ gives Eq.~(20). Differentiating with respect to $h_1$,
\begin{equation}
  \frac{\partial Z(\sigma_0)}{\partial h_1}\bigg|_{h=0}
  =2^{q-1}\tanh h_1\big|_{h_1=0}+\sigma_0(e^{a}-1)=\sigma_0(e^{a}-1),
\end{equation}
while $Z(\pm1)|_{h=0}=2^{q-1}+e^{a}-1$ for both signs, so by
Eq.~\eqref{eq:updef}
\begin{equation}
  u'=\frac{e^{a}-1}{2^{q-1}+e^{a}-1},
\end{equation}
which is Eq.~(21). At $q=2$ this is $(e^{a}-1)/(e^{a}+1)=\tanh(a/2)$: a
simplicial pair is an ordinary bond of half the coupling, as
$\delta_{\sigma_0\sigma_1}=(1+\sigma_0\sigma_1)/2$ requires. Differentiating
instead with respect to a field common to all $q-1$ other members sums $q-1$
identical derivatives and so multiplies Eq.~(21) by $q-1$; that is the
combination which enters a symmetric fixed point directly, and using it in
Eq.~(8) --- which supplies the multiplicity itself --- counts $q-1$ twice.

\subsection{The infinite-connectivity limit and its residual, Eq.~(22)}

Take one layer of cardinality $q$ at fixed chy-degree $k$, so
$\ave{\kbar}=k-1$, and impose the normalisation $\rho_qJ_q=1$ of
Ref.~[SLG26], which for $k$ complexes per vertex is $J_q=q/k$. The spinodal is
$(k-1)(q-1)u'=1$ with $a=q/(kT)$. Expanding in $\varepsilon=a$, which is $O(1/k)$,
\begin{equation}
  u'=\frac{\varepsilon+\tfrac12\varepsilon^{2}+O(\varepsilon^{3})}
          {2^{q-1}+\varepsilon+O(\varepsilon^{2})}
    =\frac{\varepsilon}{2^{q-1}}
     \Big[1+\varepsilon\big(\tfrac12-2^{1-q}\big)+O(\varepsilon^{2})\Big].
\end{equation}
Substituting and writing $T^{\ast}=q(q-1)/2^{\,q-1}$,
\begin{equation}
  \frac{T^{\ast}}{T}\Big(1-\frac1k\Big)
  \Big[1+\frac{q}{kT}\big(\tfrac12-2^{1-q}\big)\Big]=1 ,
\end{equation}
so that at leading order
\begin{equation}
  T=T^{\ast}-\frac{C_q}{k}+O(k^{-2}),
  \qquad
  C_q=T^{\ast}-q\big(\tfrac12-2^{1-q}\big)
     =\frac{q\big(2q-2^{\,q-1}\big)}{2^{\,q}} .
  \label{eq:residual}
\end{equation}
Two things follow, and the second is the one the main text states without
deriving. The limit is $T^{\ast}=q(q-1)/2^{q-1}$, which is Eq.~(8) of
Ref.~[SLG26]. And $C_q$ vanishes precisely when $2q=2^{\,q-1}$, that is at
$q=4$ and nowhere else: the leading finite-connectivity correction is absent at
the cardinality which is tricritical in the mean-field theory, and $T^\ast$ is
approached there as $1/k^{2}$ rather than $1/k$. At $q=3$, $C_3=3/4$, which
gives $1.5\times10^{-2}$, $3.8\times10^{-3}$ and $9.4\times10^{-4}$ at $k=50$,
$200$ and $800$ --- the residuals quoted in the main text, now with the
coefficient behind them. Table~\ref{tab:resid} runs the check across
cardinality.

\begin{table}[t]
\caption{\label{tab:resid}Equation~\eqref{eq:residual} against the spinodal
computed from Eq.~(8). $C_q$ changes sign between $q=4$ and $q=5$, and vanishes
at $q=4$.}
\begin{ruledtabular}
\begin{tabular}{rrrrrr}
$q$ & $T^{\ast}$ & $C_q$ & $T-T^{\ast}$ at $k=200$ & $-C_q/200$ & $T-T^{\ast}$ at $k=800$ \\
\colrule
2 & 1.00000 & $ 1.0000$ & $-5.008\times10^{-3}$ & $-5.000\times10^{-3}$ & $-1.251\times10^{-3}$ \\
3 & 1.50000 & $ 0.7500$ & $-3.763\times10^{-3}$ & $-3.750\times10^{-3}$ & $-9.383\times10^{-4}$ \\
4 & 1.50000 & $ 0$      & $-2.222\times10^{-5}$ & $0$                   & $-1.389\times10^{-6}$ \\
5 & 1.25000 & $-0.9375$ & $+4.646\times10^{-3}$ & $+4.688\times10^{-3}$ & $+1.169\times10^{-3}$ \\
6 & 0.93750 & $-1.8750$ & $+9.296\times10^{-3}$ & $+9.375\times10^{-3}$ & $+2.339\times10^{-3}$ \\
8 & 0.43750 & $-3.5000$ & $+1.721\times10^{-2}$ & $+1.750\times10^{-2}$ & $+4.356\times10^{-3}$ \\
\end{tabular}
\end{ruledtabular}
\end{table}

\section{Hitting set}
\label{sec:hs}

\subsection{The regular solution, Eq.~(32)}

For $L$ complexes of cardinality $K$ at every node, symmetry makes every
message equal and Eq.~(31) closes on the pair $(h,v)$,
\begin{equation}
  h=-\mu+(L-1)v,\qquad
  v=-\ln\Big[1-\big(1+e^{h}\big)^{-(K-1)}\Big].
\end{equation}
As $\mu\to\infty$, $h\to-\infty$ and
$(1+e^{h})^{-(K-1)}=1-(K-1)e^{h}+O(e^{2h})$, so the complex step becomes
\begin{equation}
  v=-\ln\big[(K-1)e^{h}\big]=-\ln(K-1)-h .
  \label{eq:vstep}
\end{equation}
Substituting into the chy-degree step, $h=-\mu-(L-1)\ln(K-1)-(L-1)h$, that is
$Lh=-\mu-(L-1)\ln(K-1)$, and
\begin{equation}
  h_{\mathrm{RS}}=-\frac{\mu}{L}-\frac{L-1}{L}\ln(K-1),
\end{equation}
the first line of Eq.~(32). The cavity field carries a \emph{fraction} $1/L$ of
$\mu$, which is exactly what an ansatz restricted to integer multiples of $\mu$
cannot represent, and is the reason Sec.~V~D fails above cardinality two. The
full field is what the density is read from, and by
Eq.~\eqref{eq:vstep} it is finite,
\begin{equation}
  H=h+v=-\ln(K-1),\qquad
  \rho=\frac{1}{1+e^{-H}}=\frac{1}{1+(K-1)}=\frac1K ,
\end{equation}
independent of both $\mu$ and $L$: one member of every complex is taken and none
is wasted, so the counting bound is saturated.

\subsection{The entropy, Eq.~(34)}

Equation~(33) evaluated on that solution has $n=L/K$ complexes per node and
$k=L$. From $H=-\ln(K-1)$,
\begin{equation}
  Z_i=1+e^{H}=\frac{K}{K-1},
  \qquad
  Z_a=\big(1+e^{h}\big)^{K}-1=Ke^{h}+O(e^{2h}),
\end{equation}
so $\ln Z_a=\ln K+h$ as $\mu\to\infty$. Then
\begin{equation}
  \begin{split}
  s&=\frac{L}{K}\big(\ln K+h_{\mathrm{RS}}\big)
     +(1-L)\ln\frac{K}{K-1}+\frac{\mu}{K}\\
   &=\frac{L}{K}\ln K-\frac{\mu}{K}-\frac{L-1}{K}\ln(K-1)
     +(1-L)\ln K-(1-L)\ln(K-1)+\frac{\mu}{K}\\
   &=\frac{1}{K}\Big[(L+K-LK)\ln K+(L-1)(K-1)\ln(K-1)\Big],
  \end{split}
\end{equation}
the two $O(\mu)$ pieces cancelling exactly --- which is why the regular case is
exact where a heterogeneous ensemble needs the estimator time-averaged. With
$a=(L-1)(K-1)$ one has $L+K-LK=1-a$, so
\begin{equation}
  s=\frac{1}{K}\Big[a\ln(K-1)-(a-1)\ln K\Big],
\end{equation}
which is Eq.~(34).

\subsection{The hard-field threshold, Eq.~(27)}

At fixed cardinality $c$ with Poisson chy-degree, $\bar\Phi(x)=e^{\ave{k}(x-1)}$
and Eq.~(26) closes on one scalar, $\sigma=e^{-\ave{k}\sigma^{c-1}}$.
Differentiating and using the fixed-point relation,
\begin{equation}
  |F'(\sigma)|=\ave{k}(c-1)\sigma^{c-2}e^{-\ave{k}\sigma^{c-1}}
             =\ave{k}(c-1)\sigma^{c-1},
\end{equation}
so with $y=\ave{k}\sigma^{c-1}$ the fixed point reads $\sigma=e^{-y}$,
$y=\ave{k}e^{-(c-1)y}$ and the instability $|F'|=1$ is $(c-1)y=1$. Eliminating
$y=1/(c-1)$ from the second,
\begin{equation}
  \ave{k}=y\,e^{(c-1)y}=\frac{e}{c-1},
  \qquad\text{that is}\qquad \ave{k}(c-1)=e ,
\end{equation}
Eq.~(27). Counted in neighbours the point does not move with cardinality at all,
and $c=2$ returns the Bauer--Golinelli value.

\section{Core percolation}
\label{sec:core}

\subsection{The intra-complex step, Eq.~(36)}

Run leaf removal on the cavity branch. A node reached along an inclusion is
leaf-ready ($L$, probability $\lambda$), deleted ($D$, probability $\delta$) or
core-side; a complex is traversed by entering through one member and looking at
the other $c-1$. Two events matter.

A complex leaves the entering node free to be a leaf only when it offers nothing
at all, that is when every one of the other $c-1$ members is already deleted:
$\zeta=\delta^{\,c-1}$. A complex forces the entering node out only when exactly
one other member is live and leaf-ready --- with two or more live members the
entering node is not a leaf of that complex and the complex decides nothing ---
which happens in $c-1$ ways: $\kappa=(c-1)\delta^{\,c-2}\lambda$. Those are the
two lines of Eq.~(36), and at $c=2$ they degenerate to $\zeta=\delta$,
$\kappa=\lambda$, returning the graph recursion.

\subsection{No core-free branch above cardinality two}

A core-free branch is $\gamma=0$, that is $\lambda+\delta=1$. Substituting
Eq.~(35), $\bar\Phi^{(m)}(\zeta)+1-\bar\Phi^{(m)}(1-\kappa)=1$, and since
$\bar\Phi$ is a power series with non-negative coefficients and hence strictly
increasing, this requires $\zeta=1-\kappa$ in every layer. With
$\lambda=1-\delta$ that is
\begin{equation}
  f(\delta)\equiv(c-1)\delta^{\,c-2}-(c-2)\delta^{\,c-1}-1=0 .
\end{equation}
The polynomial factors,
\begin{equation}
  f(\delta)=-(1-\delta)^{2}\sum_{j=0}^{c-3}(j+1)\,\delta^{\,j},
  \label{eq:factored}
\end{equation}
which is the whole argument. At $c=2$ the sum is empty and $f\equiv0$: the
condition is an identity and the core-free branch exists at every chy-degree
below threshold. At $c\geq3$ every coefficient of the sum is strictly positive,
so $f(\delta)<0$ on $[0,1)$ and $f(1)=0$; the only root is $\delta=1$, forcing
$\lambda=0$. A chygraph carrying any layer of cardinality three or more
therefore has no core-free branch, which is why leaf removal never certifies a
cover there, and it is a property of Eq.~(36) rather than an assumption imported
from elsewhere.

The strong reading is false and worth blocking again here: this does \emph{not}
say a complex survives leaf removal whole. Leaf removal deletes a leaf's
\emph{neighbour}, and a triangle carrying one pendant edge per vertex has an
empty core. What Eq.~\eqref{eq:factored} gives is only that the core cannot
vanish identically; it can be $1.8\times10^{-4}$.

\section{Generalised belief propagation on the region graph}
\label{sec:gbp}

Section~VII of the main text closes the complex family under intersection and
assigns M\"obius counting numbers, and stops there: the counting says where the
Bethe treatment miscounts without repairing the messages, and Table~III evaluates
it on \emph{isolated} regions, which is an estimate and not a fixed point. This
section runs the messages. The algorithm is the parent-to-child form of
generalised belief propagation [YFW05], implemented in \code{gbp.GBP}.

A message lives on a directed edge of the region graph, from a region $P$ to a
direct child $R$, and is a function of the child's variables. Writing
$U_R=\{R\}\cup\mathrm{descendants}(R)$ and $f_R$ for the product of every factor
whose scope lies inside $R$, the belief on a region is
\begin{equation}
  b_R(x_R)\;\propto\;f_R(x_R)\!\!\prod_{(I\to J):\,J\in U_R,\;I\notin U_R}\!\!
  m_{I\to J}(x_J),
\end{equation}
and imposing $\sum_{x_P\setminus x_R}b_P=b_R$ cancels every message common to
the two sides and leaves
\begin{equation}
  m_{P\to R}(x_R)=\frac{\displaystyle\sum_{x_P\setminus x_R}f_{P\setminus R}(x_P)
    \prod_{N(P,R)}m}{\displaystyle\prod_{D(P,R)\setminus(P,R)}m},
  \qquad
  \begin{aligned}
    N(P,R)&=\{(I,J):J\in U_P\setminus U_R,\ I\notin U_P\},\\
    D(P,R)&=\{(I,J):J\in U_R,\ I\in U_P\setminus U_R\}.
  \end{aligned}
  \label{eq:p2c}
\end{equation}
On the two-layer region graph of a pairwise model --- one region per interaction
plus one per variable --- $D(P,R)$ holds only the edge being updated, $N(P,R)$
is the set of incoming messages to the interaction's other variables, and
Eq.~\eqref{eq:p2c} is ordinary belief propagation. So generalised belief
propagation contains belief propagation exactly as the M\"obius counting
contains the Bethe counting, and \code{test\_gbp} checks that
\code{gbp.GBP} run on that region graph reproduces an independently written
belief propagation to $10^{-10}$ on a chain, a star, a four-cycle and a theta
graph --- the loopy error included, so the two agree on the mistake as well as
on the answer.

Two conditions are enforced before anything is iterated. The counting numbers
must cover every variable once, which \code{region.RegionGraph} already checks,
and they must cover every \emph{factor} once. The second is not automatic and is
exactly what fails for the Bethe treatment of overlapping complexes: two
triangles sharing an edge count that bond in both complexes, so
$\sum_{R\supseteq\{i,j\}}c_R=2$, and \code{gbp.GBP} refuses to run on it rather
than returning a number. Under the M\"obius counting the same sum is
$1+1-1=1$.

\subsection{The two-triangle example}

On the manuscript's own example the region family closes on
$\{0,1,2\}$, $\{1,2,3\}$ and $\{1,2\}$ once the zero-counting singletons are
dropped, and that is a junction tree; generalised belief propagation on a
junction tree is exact. That is the last column of Table~III of the main text:
the error in $\ln Z$ is at the level of double precision at every coupling,
against $1.4\times10^{-3}$ to $6.6\times10^{-2}$ for the M\"obius counting
evaluated on isolated regions and $1.8\times10^{-2}$ to $1.3$ for the Bethe
counting. So the whole of what that table reports as left on the table is
recovered, and the overshoot of the M\"obius column is a property of evaluating
a counting on isolated regions rather than of the counting itself.


\subsection{What it does and does not settle}

Three limits are worth stating plainly, because the calculation is easy to
over-read.

It is exact when the closed complex family is a junction tree, which for
maximal cliques means the graph is chordal. On $60$ maximal-clique region graphs
of hyperbolic random graphs small enough to enumerate exactly --- $n=14$, $18$
and $20$ at $\ave{k}=3.5$, $5$ and $7$, ten seeds each at $\beta J=0.3$ and
$0.8$, with only the non-treelike instances kept --- every one of the $40$
chordal runs comes out exact, the worst error being $3\times10^{-12}$.
The $20$ non-chordal runs divide. Nine of them converge, at damping up to
$0.995$, and there the error runs from $1.4\times10^{-9}$ to $5.4\times10^{-3}$
against $0.48$--$1.1$ for the M\"obius counting and $0.32$--$22$ for the Bethe
counting on the same instances --- two to eight orders of magnitude better ---
with $\sum_{x_P\setminus x_R}b_P=b_R$ satisfied to $9\times10^{-8}$, so what is
left is the approximation's error and not the solver's. The other eleven do not
converge even at damping $0.999$; their residual says so and no number from them
should be quoted. It is that instability, rather than the accuracy where it
converges, that an ensemble treatment would inherit. \code{probe/gbp\_cliques.py}
runs the comparison.

It is not variationally bounded, and on a region graph built from complexes that
contain one another it can be worse than the static estimate: covering $K_4$ by
its four triangles rather than by $K_4$ itself gives an error of $10^{-2}$
against the static $4\times10^{-3}$. That case does not arise for maximal
cliques, which never nest.

And it is an \emph{instance-level} calculation. It takes an explicit list of
complexes and an explicit interaction; everything else in the manuscript works
at the level of an ensemble, where a message is a distribution over a
chy-degree class rather than a function on one region. Lifting
Eq.~\eqref{eq:p2c} to ensembles --- messages indexed by region \emph{type}, with
the intersection structure itself random --- is what would close the overlap
deficit of Sec.~VII on hyperbolic random graphs, and it remains open.

\section{Reproducing the tables and figures}
\label{sec:repro}

\begin{table}[ht]
\caption{\label{tab:repro}What produces each display item.}
\begin{ruledtabular}
\begin{tabular}{ll}
item & script \\
\colrule
Table~I   & \code{examples/clustering\_lowers\_tc.py} \\
Table~II  & \code{examples/softfield\_vs\_hardfield.py} \\
Table~III & \code{examples/gbp\_two\_triangles.py} \\
Tables~IV, V & \code{probe/prediction4.py} \\
Table~VI  & \code{probe/clique\_moments.py} \\
Fig.~1 & \code{main.tex} (Ti\emph{k}Z) \\
Fig.~2 & \code{figures.py:clustering\_panel} \\
Fig.~3 & \code{figures.py:simplicial\_panel} \\
Fig.~4 & \code{figures.py:hittingset\_panel}, from
         \code{probe/hittingset\_density.py} \\
Fig.~5 & \code{figures.py:main}, from \code{probe/prediction4.py} \\
Eq.~(18) & \code{probe/ising\_mc.py} \\
Table~S\ref{tab:resid} & \code{tests/test\_derivations.py} \\
Sec.~\ref{sec:gbp} & \code{probe/gbp\_cliques.py}, \code{tests/test\_gbp.py} \\
\end{tabular}
\end{ruledtabular}
\end{table}

The test suite is \code{pytest tests}. It covers the derivations of
Secs.~\ref{sec:branching}--\ref{sec:gbp} as well as the outputs, which is the
practice the review established and the reason this document exists.

\end{document}
